% Conceptual adaptation of ME 537 HW03 transform/inverse and frame-chain work.
% Shared layout adapted from the instructor's ME 537 homework sources.
\documentclass[12pt,letterpaper]{article}
\usepackage{mathptmx}
\usepackage[T1]{fontenc}
\usepackage{xcolor,amsmath,amssymb,xurl,listings,enumitem,tikz}
\textwidth=6.5in
\textheight=9in
\oddsidemargin=0in
\evensidemargin=0in
\topmargin=0in
\headheight=0in
\headsep=0in
\parskip=0.08in
\parindent=0in
\setlist[enumerate]{itemsep=0.04in,topsep=0.04in}
\lstset{basicstyle=\ttfamily\footnotesize,breaklines=true,columns=fullflexible,
keepspaces=true,showstringspaces=false,aboveskip=0.04in,belowskip=0.04in}
\newcommand{\R}[2]{{}^{#1}\!R_{#2}}
\newcommand{\T}[2]{{}^{#1}\!T_{#2}}
\newcommand{\courseheader}[2]{%
\setlength{\fboxrule}{2pt}\setlength{\fboxsep}{2mm}%
\fcolorbox{black}{white}{\parbox{\dimexpr\linewidth-2\fboxsep-2\fboxrule\relax}{\centering\textbf{Introduction to Robotics Homework \##1 -- Draft}}}\par
\textbf{#2}\par}
\newcommand{\conventions}{Use right-handed rotations, column vectors, radians,
and meters. $\R{A}{B}$ maps coordinates from frame $B$ into frame $A$;
$\T{A}{B}$ also includes the origin of $B$ expressed in $A$.}
\newcommand{\submission}{\textbf{Submit:} labeled hand derivations and sketches,
your changed Python source files, the requested numerical/visual evidence,
and an AI-use statement (including ``none'' if applicable). Show your reasoning
before using code to check hand calculations. Public checks are intentionally
limited. The instructor supplies the deadline and submission destination.}

\begin{document}
\courseheader{3}{Homogeneous transforms and camera-to-arm coordinates}
\textbf{Schedule:} week 3, approximately one week. Build on HW02. \conventions

\begin{enumerate}
\item \textbf{Implement rigid transforms and their inverse.}
Complete HW03's \texttt{se3} and \texttt{inv} in
\path{src/intro_robotics/transforms.py}.
\begin{enumerate}
\item Write the block form of $\T{A}{B}$ in terms of $\R{A}{B}$ and
$p^A_B$. Derive its inverse by requiring both products to be identity.
\item Implement the block operations directly. Use rotation structure rather
than a general matrix inverse. The supplied signatures permit identity
rotation and zero translation defaults; preserve those conveniences.
\item Run \texttt{python -m homework.hw03.hw03\_transforms} from the
container shell. Add a nontrivial round-trip check of your own, with both
rotation and translation nonzero. Explain its frame directions.
\end{enumerate}

\item \textbf{Locate a conveyor object in the robot base frame.}
Let $W$ be the world, $B$ the robot base, and $C$ a fixed camera.
The camera looks down at a plane $z_W=0$. You are given:
\[
\T{W}{B}=\begin{bmatrix}
0&-1&0&0.1\\1&0&0&0.1\\0&0&1&0\\0&0&0&1
\end{bmatrix},\qquad
\T{W}{C}=\begin{bmatrix}
1&0&0&0.4\\0&-1&0&-0.2\\0&0&-1&0.8\\0&0&0&1
\end{bmatrix}.
\]
A measured point is $p_C=[0.1,\ 0.05,\ 0.8]^T$ meters in camera coordinates.
This is already a 3D point; pixel-to-3D estimation belongs in later vision work.
\begin{enumerate}
\item Sketch the origins and describe the camera's three axes in world coordinates.
\item Derive and evaluate $\T{B}{C}$ by hand. State what its translation means.
\item Calculate $p_W$ and $p_B$ by hand. Check that the point lies on the
world conveyor plane. Verify by mapping $p_B$ back through $B$ into $W$.
\item For a direction $v_C=[0,1,0]^T$, calculate $v_W$ and $v_B$.
Explain why the transform's translation does not change a direction.
Use homogeneous last entries 1 for a point and 0 for a direction.
\end{enumerate}
\end{enumerate}

\newpage
\begin{enumerate}[start=3]
\item \textbf{Check that rotation and translation do not commute.}
Define $R=\operatorname{Rot}_z(\pi/2)$ and $d=[0.3,0,0]^T$ meters. Let
\[
T_R=\begin{bmatrix}R&0\\0&1\end{bmatrix},\qquad
T_d=\begin{bmatrix}I&d\\0&1\end{bmatrix}.
\]
Here the bottom row notation represents $[0,0,0,1]$.
\begin{enumerate}
\item Find $T_R T_d$ and $T_d T_R$ by hand. Apply each product to the
same point $[0.2,0,0]^T$ and sketch the two resulting positions.
\item Explain the geometric difference. For each product, is simply negating
the original translation sufficient to invert it? Use your block inverse to explain.
\end{enumerate}

\item \textbf{Check the frame chain with ROS TF.}
ROS keeps a tree of named frames called TF. Each entry is the pose of a
\emph{child} frame expressed in its \emph{parent} frame: a matrix
$\T{\text{parent}}{\text{child}}$, exactly as in this assignment.
\begin{enumerate}
\item Write $\T{C}{O}$ by hand for an object frame $O$ located at $p_C$ and
oriented like the world. Then complete \texttt{frame\_transforms()} in
\path{homework/hw03/frame_broadcaster.py}: return the poses of
\texttt{base} and \texttt{camera} in \texttt{world}, and of \texttt{object}
in \texttt{camera}, built with your \texttt{se3} and \texttt{inv}.
\item Publish the frames, then query TF from a second container shell:
\begin{lstlisting}
python3 -m homework.hw03.frame_broadcaster
ros2 run tf2_ros tf2_echo base object
ros2 run tf2_ros tf2_echo base camera
\end{lstlisting}
Compare the printed translations and matrices with your hand answers for
$p_B$ and $\T{B}{C}$ in problem 2. TF found these by composing the tree for
you. Write the product of your published matrices (and any inverse) that
gives \texttt{base} to \texttt{object}.
\item In RViz, set the Fixed Frame to \texttt{world} and add a TF display.
Save a screenshot showing all four frames. Which matrix would you publish if
\texttt{camera}'s parent were \texttt{base}? Where would RViz draw the camera
if you published $\T{B}{C}$ with parent \texttt{world}, and why?
\end{enumerate}
\end{enumerate}

\textbf{What to understand:} transform direction and composition order,
translation expressed in the reference frame, rigid-transform inversion,
and the difference between points and directions. A TF tree stores exactly
these parent-to-child transforms. No DH convention is used.

\submission
\end{document}
